% E15 paragraph; numbers from results/final/E15_scalability.json and E15_optimality_gap.json (paper/e15_numbers.py).
% Plan: request latency on the real registry and on cloned registries; engine build; exact vs greedy.
\emph{Cost and scalability (E15).} On the real registry a request takes 49\,ms at the median and 191\,ms at the
95th percentile on two virtual CPUs; eligibility takes 0.08\,ms and the per-goal retrieval, timed separately,
45\,ms, against 12\,ms for one dense ranking of the request (RAG-MCP). On registries of cloned contracts, which are
near-copies of the real tools, the median grows to 58\,ms at 1,000 tools, 281\,ms at 10,000 and 2.7\,s at 100,000
(16.1\,s at the 95th percentile), almost all of it retrieval (RAG-MCP: 206\,ms). The bitmap test alone takes
0.024\,ms over $10^6$ contracts, an index microbenchmark, not the full engine at that size.
Building the engine takes 2.2\,s, 2.8\,s, 30.3\,s and 2,224\,s: the implementation materialises the pairwise risk
relation eagerly, which is quadratic in the number of tools. Compositions use the exact search for 26--28 of the 30
tasks at every size; abstentions ($\chi_1$, and one $\chi_5$) rise from 2 tasks to 9 at 10,000 and 10 at 100,000
tools. On generated cover instances, measured with the objective the search minimises (Eq.~(6)), the greedy regime
is optimal in 7, 8 and 2 of 10 instances with 1, 3 and 5 goals and on average 0.7\%, 2.3\% and 14.5\% above the
optimum (at most 3.0\%, 19.8\% and 38.1\%), while the exact search takes 3\,ms, 21\,ms and 1.1\,s. With 8 and 12 goals
the exact search exhausts its 20,000-node budget in all ten instances (after 12.0\,s and 18.2\,s) and falls back to
the greedy regime (27\,ms and 69\,ms), which carries no optimality guarantee; the gap could not be measured there.
